% Generated by report/tools/tables.py from the result files of the repository. Do not edit.
\begin{table}[tbp]
\centering
\footnotesize
\setlength{\tabcolsep}{3.2pt}
\begin{threeparttable}
\caption[GRPO : modèles, récompenses, graines]{Volet GRPO (\file{sim/llm/LLM\_PROTOCOL.md}) : deux petits modèles ouverts, Qwen2.5 à 0,5 milliard de paramètres et Llama 3.2 à 1 milliard, entraînés 100 pas sur la tâche \texttt{mul2} (produit de deux nombres de 12 à 99), avec la vraie récompense puis une récompense tirée à pile ou face, cinq graines (0 à 4). \textbf{En gras}, les intervalles qui excluent zéro. Sources : \file{sim/llm/results/summary\_5.json}, \file{gains\_5.json}, \file{extra\_5.json}.}
\label{tab:grpo}
\begin{tabular}{@{}lllrrrrrr@{}}
\toprule
 & & & & \multicolumn{3}{c}{Réussite au premier essai (\%)} & \multicolumn{2}{c}{Réponses coupées (\%)} \\
\cmidrule(lr){5-7}\cmidrule(lr){8-9}
Modèle & Récompense & Graine & Matériel & pas 0 & pas 100 & gain (pts) & pas 0 & pas 100 \\
\midrule
Qwen2.5 0,5B & vraie & 0 & GPU (MPS) & 10 & 36 & +26 & 77 & 29 \\
 &  & 1 & GPU (MPS) & 10 & 61 & +51 & 75 & 0 \\
 &  & 2 & CPU & 11 & 13 & +2 & 76 & 0 \\
 &  & 3 & CPU & 10 & 13 & +4 & 74 & 0 \\
 &  & 4 & CPU & 11 & 0 & −11 & 73 & 0 \\
\addlinespace
 & au hasard & 0 & GPU (MPS) & 10 & 1 & −9 & 77 & 59 \\
 &  & 1 & CPU & 12 & 0 & −11 & 71 & 88 \\
 &  & 2 & CPU & 11 & 0 & −11 & 76 & 98 \\
 &  & 3 & CPU & 10 & 0 & −10 & 74 & 16 \\
 &  & 4 & CPU & 11 & 0 & −11 & 73 & 86 \\
\addlinespace
Llama 3.2 1B & vraie & 0 & GPU (MPS) & 25 & 20 & −4 & 2 & 9 \\
 &  & 1 & GPU (MPS) & 25 & 48 & +22 & 1 & 0 \\
 &  & 2 & GPU (MPS) & 23 & 25 & +2 & 2 & 3 \\
 &  & 3 & GPU (MPS) & 25 & 27 & +2 & 3 & 1 \\
 &  & 4 & GPU (MPS) & 23 & 33 & +10 & 3 & 1 \\
\addlinespace
 & au hasard & 0 & GPU (MPS) & 25 & 0 & −25 & 2 & 99 \\
 &  & 1 & GPU (MPS) & 25 & 1 & −24 & 1 & 100 \\
 &  & 2 & GPU (MPS) & 23 & 0 & −23 & 2 & 2 \\
 &  & 3 & GPU (MPS) & 24 & 0 & −24 & 2 & 16 \\
 &  & 4 & GPU (MPS) & \multicolumn{5}{l}{échec deux fois} \\
\midrule
\multicolumn{7}{@{}l}{Gain moyen en points [IC bootstrap 95~\%]} & \multicolumn{2}{r}{} \\
\multicolumn{7}{@{}l}{\quad Qwen2.5 0,5B, vraie, 5 graines} & \multicolumn{2}{r}{+14 [−3\,;\,+34]} \\
\multicolumn{7}{@{}l}{\quad Qwen2.5 0,5B, au hasard, 5 graines} & \multicolumn{2}{r}{\textbf{−10 [−11\,;\,−9]}} \\
\multicolumn{7}{@{}l}{\quad Llama 3.2 1B, vraie, 5 graines} & \multicolumn{2}{r}{+6 [−0,4\,;\,+15]} \\
\multicolumn{7}{@{}l}{\quad Llama 3.2 1B, au hasard, 4 graines} & \multicolumn{2}{r}{\textbf{−24 [−24\,;\,−23]}} \\
\multicolumn{7}{@{}l}{\quad Qwen2.5 0,5B, vraie, CPU seul, graines 2, 3, 4} & \multicolumn{2}{r}{−2 [−11\,;\,+4]} \\
\multicolumn{7}{@{}l}{\quad Qwen2.5 0,5B, au hasard, CPU seul, graines 1, 2, 3, 4} & \multicolumn{2}{r}{\textbf{−11 [−11\,;\,−10]}} \\
\multicolumn{7}{@{}l}{\quad Llama 3.2 1B, au hasard, graine 4 = plus petit gain, −25 } & \multicolumn{2}{r}{−24 [−24\,;\,−23]} \\
\multicolumn{7}{@{}l}{\quad Llama 3.2 1B, au hasard, graine 4 = plus grand gain, +51 } & \multicolumn{2}{r}{−9 [−24\,;\,+21]} \\
\bottomrule
\end{tabular}
\begin{tablenotes}[flushleft]\scriptsize
\item[] Réussite au premier essai : pass@1 échantillonné sur les problèmes de test, en pourcentage ; gain : pas 100 moins pas 0, en points ; réponses coupées : part des réponses qui atteignent la limite de 160 jetons sans donner de résultat, toutes comptées fausses. Matériel : GPU Apple (MPS) ou CPU ; trois essais Qwen sont passés sur CPU après des gradients non finis sur MPS, si bien que matériel et graine se confondent pour Qwen. Partie basse : statistique préregistrée, moyenne des gains par graine et intervalle bootstrap percentile à 95~\% (10\,000 tirages, \texttt{random.Random(0)}), très grossier avec cinq graines. La configuration Llama, récompense au hasard, est résumée sur ses quatre essais finis, avec la sensibilité au gain manquant. Part coupée au pas 0 : première évaluation de chaque fichier \file{sim/llm/results/<modèle>/mul2/<récompense>/seed<g>.jsonl}.
\end{tablenotes}
\end{threeparttable}
\end{table}
